\section{Nash 均衡存在但可能不唯一也不好}
\label{sec:17-Constrained-Guarantees/06-Games-and-Equilibria/02}

两家外卖平台打了一年补贴战。双方都清楚，如果一起把补贴降下来，各自每季度能多留下三个亿；这话在行业会议上被说过不止一次。可谁也没有先降，一年下来两家各自只剩一个亿。

事后复盘常把这归结为``缺乏互信''或者``短视''。把数字摆开就会看到不是这么回事。设双方各有两个选择，把季度利润写成一对数，前一个是平台甲的，后一个是平台乙的：

\[
\begin{array}{c|cc}
 & \text{乙低补贴} & \text{乙高补贴}\\\hline
\text{甲低补贴} & (3,3) & (0,5)\\
\text{甲高补贴} & (5,0) & (1,1)
\end{array}
\]

现在站在甲的位置上算。若乙选低补贴，甲选低得 \(3\)、选高得 \(5\)，该选高。若乙选高补贴，甲选低得 \(0\)、选高得 \(1\)，还是该选高。无论乙做什么，甲选高补贴都严格更好。乙的算式与甲对称，结论相同。

于是 \((\text{高},\text{高})\) 是唯一一个双方都无法单方面改善的组合，落在那里各得 \(1\)；而 \((\text{低},\text{低})\) 给出 \((3,3)\)，对两家都更好，却没有一方愿意留在那里。不需要短视，不需要不信任，只需要各自算清自己那本账。

\subsection{均衡是最优反应映射的不动点}

把上面的推理写成一般形式。\(n\) 个参与者，第 \(i\) 个的策略集是 \(S_i\)，效用是 \(u_i(s_1,\dots,s_n)\)。记 \(s_{-i}\) 为除 \(i\) 之外所有人的策略。一个策略组合 \(s^\star\) 叫 Nash 均衡，如果对每个 \(i\) 和每个 \(s_i\in S_i\) 都有

\[
u_i(s_i^\star,s_{-i}^\star)\;\ge\;u_i(s_i,s_{-i}^\star).
\]

这句话只说了一件事：在别人不动的前提下，谁也无法靠自己变好。它不说这个组合好，不说它公平，也不说博弈一定会走到那里。

存在性是 Nash 在 1950 年给出的。定义每个人的最优反应对应

\[
\mathrm{BR}_i(s)=\argmax_{s_i\in S_i}u_i(s_i,s_{-i}),
\qquad
\mathrm{BR}(s)=\mathrm{BR}_1(s)\times\cdots\times\mathrm{BR}_n(s).
\]

按定义，\(s^\star\) 是 Nash 均衡当且仅当 \(s^\star\in\mathrm{BR}(s^\star)\)。均衡就是最优反应这个映射的不动点，剩下的全部工作是找一条定理保证不动点存在。

证明骨架用 Kakutani 不动点定理，它要求四件事：定义域是欧氏空间中的非空凸紧集，对应的取值非空、取值是凸集、图是闭的。把纯策略换成混合策略之后逐条核对。定义域是 \(\Delta_1\times\cdots\times\Delta_n\)，有限个单纯形的乘积，凸性来自\hyperref[sec:09-Convex-Optimization/02-Convex-Sets/01]{``线段、凸组合与凸包''}里那条``按比例掺着用仍然可行''，紧性则是\hyperref[sec:12-Real-Analysis-Support/03-Topology-Basics/02]{``紧致性把无限局部信息压成有限控制''}里那个让极值真的取得到的条件。\(u_i\) 关于 \(p_i\) 是线性的，在紧集上必取到最大值，所以 \(\mathrm{BR}_i\) 非空；线性函数的最大值点集是单纯形的一个面，是凸集；效用连续保证了图闭，也就是取极限不会跑出最优反应集合。四条齐备，Kakutani 给出不动点，Nash 均衡存在。

零和情形是这个结论的特例：两人零和博弈的 Nash 均衡恰好是\hyperref[sec:17-Constrained-Guarantees/06-Games-and-Equilibria/01]{``最坏情况下的最优与交换次序''}里的鞍点，均衡收益就是博弈的值。这也解释了为什么零和博弈的均衡是可以互换的——任意两个均衡的收益都等于同一个值——而一般和博弈里没有这条性质。

条件对账。有限策略集与混合策略是紧凸性的来源，把任何一条抽掉都可能出问题。策略集不紧时最优反应可能取不到：两人各从 \([0,1)\) 中报一个数，报得大的赢，谁都想再往上挪一点，没有均衡。效用不连续时同样失效，而且这个例子更贴近现实：两家公司各自在 \([0,1]\) 上报价，报低的一方拿走全部订单，平价时对半分。这里策略集紧、凸，可效用在对角线上跳变，最优反应对应的图不闭——对手报 \(y\) 时我想报``比 \(y\) 略低一点''，而``略低一点''不是一个可以取到的策略。均衡不存在。把价格离散成整数分，均衡立刻回来，因为策略集重新变成有限的。

\subsection{Kakutani 给存在性，不给别的}

不动点在本书里出现过好几次，把它们并排看能看清这次得到的是什么、没得到什么。\hyperref[sec:18-Synthesis/02-Recurring-Ideas/03]{``迭代：把解写成一列近似的极限''}讲的是这个动作的一般形态：想求的东西写成 \(x=T(x)\)，然后反复作用 \(T\)。

\hyperref[sec:08-Numerical-Analysis/02-Root-Finding/02]{``不动点与 Newton 法''}里的 Banach 压缩映射定理给了一整套东西：不动点存在、唯一、迭代必然收敛、收敛速度是几何式的而且比率写得出来。代价是一个很强的前提——\(T\) 必须是压缩的。\hyperref[sec:16-Control-and-Reinforcement-Learning/02-Reinforcement-Learning/02]{``Bellman 方程与值迭代、策略迭代''}里的 Bellman 算子恰好满足这个前提，折扣因子 \(\gamma\) 就是压缩系数，于是值迭代是一个真正可以跑的算法。

Kakutani 什么也不要求，也什么都不多给。它只说不动点存在。最优反应对应一般不是压缩映射——在囚徒困境里它甚至是常值的，在石头剪刀布里它绕圈——所以没有一个与值迭代对应的通用算法可以照着定义把均衡算出来。这不是懒得找算法：计算一个两人一般和博弈的 Nash 均衡是 PPAD-完全问题，而判断是否存在满足某些附加条件（比如社会总收益超过某个门槛）的均衡是 NP-难的，\hyperref[sec:09-Convex-Optimization/09-Complexity-and-Approximation/01]{``P、NP 与 NP-hard：给问题按计算难度分类''}讲的那条分界线在这里同样适用。

零和两人博弈是个例外，它的均衡可以用线性规划在多项式时间内解出来，因为极小极大定理把它变成了一个凸问题。从零和到一般和，丢掉的不只是``博弈的值''这个概念，还有可计算性。

\subsection{存在不等于唯一}

Nash 定理保证至少有一个均衡，一个字也没说只有一个。

两家公司要选一个数据交换格式，选同一个都省事，选不同的都白费力气。收益是：都选 X 各得 \(2\)，都选 Y 各得 \(2\)，选得不一样各得 \(0\)。这里有两个纯策略均衡，还有一个双方各以 \(1/2\) 概率随机的混合均衡（收益各 \(1\)，比两个纯均衡都差）。三个均衡并存，博弈论本身不告诉你会走到哪一个。

这件事的后果比看上去严重。若有唯一均衡，``理性参与者会怎么做''还算有个答案；有多个均衡时，这个问题的答案必须由博弈之外的东西来定——历史惯例、谁先公开表态、哪一个显得更``自然''、或者某种演化动力学的初始条件。均衡选择不是一个可以从收益矩阵推出来的东西，它需要额外的假设，而那些假设通常是经验判断而不是定理。

于是要留心一种常见的误用：算出一个 Nash 均衡，就宣称``系统会稳定在这里''。这句话需要三个前提，每一个都可能不成立——均衡唯一、参与者能找到它、以及某种动力学确实把系统带过去。第三条是下一个单元回来处理的事。

\subsection{存在也不等于好}

囚徒困境说的是另一件事：均衡可能对所有人都更差。

\begin{center}
\begin{tikzpicture}[x=1cm,y=1cm,font=\small,>=stealth]
  \draw[black!55] (0,0) rectangle (5.2,3.0);
  \draw[black!55] (2.6,0) -- (2.6,3.0);
  \draw[black!55] (0,1.5) -- (5.2,1.5);
  \fill[green!12] (0,1.5) rectangle (2.6,3.0);
  \fill[red!12] (2.6,0) rectangle (5.2,1.5);
  \draw[black!55] (0,0) rectangle (5.2,3.0);
  \draw[black!55] (2.6,0) -- (2.6,3.0);
  \draw[black!55] (0,1.5) -- (5.2,1.5);
  \node at (1.3,2.25) {\((3,\,3)\)};
  \node at (3.9,2.25) {\((0,\,5)\)};
  \node at (1.3,0.75) {\((5,\,0)\)};
  \node at (3.9,0.75) {\((1,\,1)\)};
  \node[font=\footnotesize,black!70] at (1.3,3.28) {乙低补贴};
  \node[font=\footnotesize,black!70] at (3.9,3.28) {乙高补贴};
  \node[font=\footnotesize,black!70,anchor=east] at (-0.15,2.25) {甲低补贴};
  \node[font=\footnotesize,black!70,anchor=east] at (-0.15,0.75) {甲高补贴};
  \draw[->,blue!65!black,thick] (1.05,1.95) -- (1.05,1.05);
  \draw[->,blue!65!black,thick] (3.65,1.95) -- (3.65,1.05);
  \draw[->,red!70!black,thick] (1.85,2.55) -- (3.35,2.55);
  \draw[->,red!70!black,thick] (1.85,0.35) -- (3.35,0.35);
  \node[anchor=west,font=\footnotesize,blue!65!black] at (5.45,1.5) {甲单方面改进};
  \node[anchor=west,font=\footnotesize,red!70!black] at (5.45,2.55) {乙单方面改进};
  \node[anchor=west,align=left,font=\footnotesize,black!75] at (5.45,0.6)
    {四个箭头汇入同一格：\\ 那格是唯一均衡，\\ 而左上角对双方都更好};
  \node[font=\scriptsize,black!60] at (2.6,-0.5) {收益记为（甲，乙），数值越大越好};
\end{tikzpicture}

\emph{每个箭头都是一次合理的单方面改进，四个箭头全部指向右下角。均衡的定义只要求没人能靠自己变好，它对``大家一起变好''这件事没有任何管辖权。}
\end{center}

图里左上角那一格在帕累托意义上支配右下角：\(3>1\) 且 \(3>1\)，两人同时更好。可它不是均衡，因为从那里出发，任何一方单方面切到高补贴都能把自己从 \(3\) 提到 \(5\)。均衡要求的是对单方面偏离稳定，而``一起变好''需要联合偏离，那不在定义的管辖范围内。

这条结论的适用面比补贴战宽得多。军备竞赛、公共渔场的过度捕捞、平台之间的算法竞价、同一个组织里各团队为抢资源而虚报需求，结构都是这个矩阵。它们的共同点不是参与者愚蠢，而是每个人的最优选择产生的外部影响没有落在自己的账本上。

值得把两种``坏''分清楚。上一小节的坏是不确定——多个均衡，不知道去哪个。这一小节的坏是确定的——唯一均衡，而且确定地差。前者需要额外的选择原则，后者需要改动博弈本身。两者都不能靠``让参与者更理性''来解决，恰恰相反，参与者越理性，右下角越牢固。

那么``更差''到底差多少？这个问题得先能被量化，才谈得上判断某个机制是不是值得改。下一节给它一个比值，并用一个反直觉的例子说明，往系统里增加选项也可能让所有人变差。
